\subsection{Information Carried by a Single Qubit}

This section addresses a fundamental question: \textbf{\emph{if a qubit can exist in infinitely many states, how much information can it actually convey?}} A general single-qubit state is given by:
\begin{equation*}
    \ket{\psi} = a\ket{0} + b\ket{1} \qquad a,b \in \mathbb{C}, \quad \left|a\right|^2 + \left|b\right|^2 = 1
\end{equation*}
Because $a$ and $b$ are complex numbers, a qubit can be prepared in \textbf{infinitely many distinct states}. This might suggest that a qubit can encode an infinite amount of information. However, this is \textbf{not} the case. When we perform a \textbf{measurement}, we obtain only a \textbf{single classical outcome} (either 0 or 1). Therefore, \hl{a single qubit can yield at most one classical bit of information when measured}. This apparent paradox, \textbf{infinite possible states but limited measurable information}, is one of the most important conceptual aspects of quantum information theory.

\highspace
\begin{flushleft}
    \textcolor{Green3}{\faIcon{question-circle} \textbf{Why can't we determine the Amplitudes $a$ and $b$ of a single qubit state?}}
\end{flushleft}
Suppose we are given a qubit in an unknown state $\ket{\psi}$. Can we determine $a$ and $b$ from a single measurement? The answer is \textbf{no}, because:
\begin{itemize}
    \item[\textcolor{Red2}{\faIcon{times}}] \textbf{Measurement is probabilistic}. A single measurement provides only one outcome (0 or 1), which is insufficient to determine the amplitudes.
    \item[\textcolor{Red2}{\faIcon{times}}] \textbf{Measurement disturbs the state}. After measurement, the qubit collapses to the observed basis state, preventing further measurements on the original state.
    \item[\textcolor{Red2}{\faIcon{times}}] \textbf{State estimation requires many copies}. To estimate $a$ and $b$, we need \textbf{quantum state tomography}, which relies on performing measurements on \textbf{many identically prepared qubits}.
\end{itemize}
This \textbf{irreversible collapse} means that the original information encoded in the amplitudes is generally lost after measurement.

\highspace
\begin{flushleft}
    \textcolor{Red2}{\faIcon{exclamation-triangle} \textbf{The No-Cloning Theorem}}
\end{flushleft}
One might think of copying the qubit before measurement to extract more information. However, quantum mechanics forbids this through the \definition{No-Cloning Theorem}.

\begin{theorem}[\textbf{No-Cloning}]\label{thm:no-cloning}
    It is impossible to create an identical copy of an arbitrary unknown quantum state.
\end{theorem}
\begin{proof}
    Assume there exists a unitary operator $U$ that can clone any state $\ket{\psi}$:
    \begin{equation*}
        U\ket{\psi}\ket{0} = \ket{\psi}\ket{\psi} \quad \forall \ket{\psi}
    \end{equation*}
    For two arbitrary states $\ket{\psi}$ and $\ket{\phi}$, linearity of quantum mechanics implies:
    \begin{equation*}
        U\left(\alpha\ket{\psi} + \beta\ket{\phi}\right)\ket{0} = \alpha\ket{\psi}\ket{\psi} + \beta\ket{\phi}\ket{\phi}
    \end{equation*}
    Which is \textbf{not equal} to:
    \begin{equation*}
        \left(\alpha\ket{\psi} + \beta\ket{\phi}\right)\left(\alpha\ket{\psi} + \beta\ket{\phi}\right) = \alpha^2\ket{\psi}\ket{\psi} + \alpha\beta\ket{\psi}\ket{\phi} + \alpha\beta\ket{\phi}\ket{\psi} + \beta^2\ket{\phi}\ket{\phi}
    \end{equation*}
    This contradiction proves that perfect cloning is impossible, since the two expressions cannot be equal for all $\ket{\psi}$ and $\ket{\phi}$.
\end{proof}

\noindent
Since we cannot clone a qubit, we cannot create multiple copies to determine its exact state from a single instance.

\begin{table}[!htp]
    \centering
    \begin{tabular}{@{} p{9em} l l @{}}
        \toprule
        Feature & Classical Bit & Qubit \\
        \midrule
        Possible states & 0 or 1 & Qubit \\[.2em]
        Info from measurement & 1 bit & At most 1 bit \\[.2em]
        Measurement disturbance & No & Yes \\[.2em]
        Ability to copy & Yes & No (No-Cloning) \\[.2em]
        State determination & Single observation & Requires many identical copies \\
        \bottomrule
    \end{tabular}
    \caption{Comparison of classical bits and qubits in terms of information content and measurement properties.}
\end{table}